% eksperymentalne tłumaczenie na włoski

%

\subsection{Środek ciężkości i trzy środkowe} % lang-pl
\subsection{Baricentro e tre mediane} % lang-it
\begin{definition}[mediana] % lang-it
\index{mediana}% % lang-it
    Dato un triangolo. % lang-it
    Il segmento che unisce uno qualsiasi dei suoi vertici con il punto medio del lato opposto si chiama mediana. % lang-it
\end{definition} % lang-it


\begin{definition}[środkowa] % lang-pl
\index{środkowa}% % lang-pl
    Dany jest pewien trójkąt. % lang-pl
    Odcinek łaczący którykolwiek z jego wierzchołków ze środkiem przeciwległego boku nazywamy środkową. % lang-pl
\end{definition} % lang-pl
Da non confondere con il segmento medio! % lang-it

Nie mylić z linią środkową! % lang-pl
Środkowe dzielą trójkąt na sześć mniejszych o równych polach. % lang-pl
Le mediane dividono il triangolo in sei triangoli più piccoli di uguale area. % lang-it
% TODO: rysunek równych sześciu pól.

% TODO: długość środkowej, Zetel s. 32

\begin{proposition}
    Niech $m_a, m_b, m_c$ oznacza długości trzech środkowych trójkąta, zaś $p$ połowę jego obwodu. % lang-pl
    
    Wtedy % lang-pl
    Siano $m_a, m_b, m_c$ le lunghezze delle tre mediane di un triangolo e $p$ il semiperimetro. % lang-it
    
    Allora % lang-it
    \begin{equation}
        \frac 32 p \le m_a + m_b + m_c \le 2p.
    \end{equation}
    % Posamentier, Alfred S., and Salkind, Charles T., Challenging Problems in Geometry, Dover, 1996: pp. 86–87.
\end{proposition}

Napisze o tym Coxeter \cite[s. 26]{coxeter_1967}. % lang-pl

Guzicki \cite[s. 357]{guzicki_2021} wytłumaczy dodatkowo, czemu żadnej ze stałych: $\frac 34$, $1$ nie można poprawić. % lang-pl
Ne scriverà Coxeter \cite[p. 26]{coxeter_1967}. % lang-it
Guzicki \cite[p. 357]{guzicki_2021} spiegherà inoltre perché nessuna delle costanti $\frac 34$, $1$ può essere migliorata. % lang-it


Mamy też podobną nierówność: % lang-pl
Abbiamo anche una disuguaglianza simile: % lang-it

\begin{proposition}
    Niech $a, b, c$ oznacza długości boków trójkąta. % lang-pl
    
    Wtedy % lang-pl
    Siano $a$, $b$, $c$ le lunghezze dei lati di un triangolo. % lang-it
    
    Allora % lang-it
    \begin{equation}
        \frac 32 \le \frac{a}{b+c} + \frac{b}{a+c} + \frac{c}{a+b} < 2.
    \end{equation}
\end{proposition}

Lewa połowa jest prawdziwa dla wszystkich dodatnich liczb $a, b, c$ i nosi nazwę nierówności (Alfreda Mortimera) Nesbitta, który znajdzie ją około 1902 roku. % lang-pl

La disuguaglianza a sinistra vale per tutti i numeri positivi $a, b, c$ ed è nota come disuguaglianza di (Alfred Mortimer) Nesbitt, che la troverà intorno al 1902. % lang-it

\begin{proposition}
\label{srodkowe_przecinaja_sie}%
\index{środek ciężkości}% % lang-pl
    Trzy środkowe trójkąta przecinają się w jednym punkcie zwanym środkiem ciężkości i dzielą się w~stosunku $2$ do $1$, licząc od wierzchołków. % lang-pl
\index{baricentro}% % lang-it
    
    Le tre mediane di un triangolo si incontrano in un unico punto, chiamato baricentro, e vengono divise da esso nel rapporto $2$ a $1$, contando dai vertici. % lang-it
\end{proposition}
% % Coxeter, Introduction to Geometry, s. 10 <- przeczytaj to, nie tylko cytuj! + ćwiczenia: 3/4 <= 1
% Neugebauer s. 54: oznaczenie G, bo gravis - ciężki
% Zetel, s. 25 z tw. Cevy

Środkowe to po angielsku \emph{medians}, przecinają się w \emph{centroid}. % lang-pl

Polska nazwa weźmie się z tego, że środek ciężkości fizycznego modelu trójkąta wykonanego z jednolitego materiału znajduje się właśnie tam. % lang-pl
Po angielsku środkowe nazywają się \emph{medians}, a ich punkt przecięcia \emph{centroid}. % lang-it

Angielska nazwa powstanie w 1814 roku z powodu niechęci do tego, co nie jest czysto geometryczne; niechęci, której nie będzie widać w innych europejskich językach. % lang-pl
Il nome polacco deriva dal fatto che il centro di gravità di un modello fisico di triangolo realizzato con materiale omogeneo si trova proprio in quel punto. % lang-it
Il termine inglese nascerà nel 1814 a causa di una certa avversione verso ciò che non è puramente geometrico; un'avversione che non si osserverà nelle altre lingue europee. % lang-it


(Środek ciężkości brzegu trójkąta nazywa się punktem Spiekera, patrz fakt \ref{punkt_spiekera}). % lang-pl
(Il centro di gravità del bordo di un triangolo si chiama punto di Spieker, vedi il fatto \ref{punkt_spiekera}). % lang-it
\index{punkt!Spiekera}%

Hartshorne \cite[s. 52-54]{hartshorne2000} wnioskuje \ref{srodkowe_przecinaja_sie} z faktu \ref{hartshorne_52} (później zaś \cite[s. 119-120]{hartshorne2000} powtórzy dowód z~maszynerią geometrii analitycznej). % lang-pl

Hartshorne \cite[p. 52-54]{hartshorne2000} deduce \ref{srodkowe_przecinaja_sie} dal fatto \ref{hartshorne_52} (più tardi \cite[p. 119-120]{hartshorne2000} ripeterà la dimostrazione usando gli strumenti della geometria analitica). % lang-it
\index[persons]{Archimedes}% % lang-pl
\index[persons]{Archimede}% % lang-it
Podobnie zrobią Pompe \cite[s. 41]{pompe_2022}, Guzicki \cite[s. 220]{guzicki_2021} albo Bogdańska, Neugebauer \cite[s. 53]{guzicki_2021}. % lang-pl

Zetel \cite[s. 14, 25]{zetel_2020} pokaże, że twierdzenie Cevy daje natychmiast dowód istnienia środka ciężkości. % lang-pl
Faranno lo stesso Pompe \cite[p. 41]{pompe_2022}, Guzicki \cite[p. 220]{guzicki_2021} e Bogdańska, Neugebauer \cite[p. 53]{guzicki_2021}. % lang-it

\index{twierdzenie!Cevy}% % lang-pl
Zetel \cite[p. 14, 25]{zetel_2020} mostrerà che il teorema di Ceva fornisce immediatamente una dimostrazione dell'esistenza del baricentro. % lang-it
\index{teorema!di Ceva}% % lang-it
% TODO: Ich długość można obliczyć z https://en.wikipedia.org/wiki/Apollonius%27s_theorem => to jest z https://en.wikipedia.org/wiki/Median_(geometry)#Formulas_involving_the_medians'_lengths
Coxeter \cite[s. 26]{coxeter_1967} pokaże, jak otrzymać ten sam wynik z (VI.2) oraz (VI.4). % lang-pl
Coxeter \cite[p. 26]{coxeter_1967} mostrerà come ottenere lo stesso risultato da (VI.2) e (VI.4). % lang-it


Historia faktów \ref{wysokosci_przecinaja_sie} oraz \ref{srodkowe_przecinaja_sie} ma wiele elementów wspólnych, więc omówimy je razem. % lang-pl
La storia dei fatti \ref{wysokosci_przecinaja_sie} e \ref{srodkowe_przecinaja_sie} presenta molti elementi in comune, perciò li discuteremo insieme. % lang-it

Żaden z nich nie pojawi się w Elementach Euklidesach i właściwie nie wiadomo, kto odkryje je pierwszy. % lang-pl
Nessuno dei due compare negli Elementi di Euclide e non si sa realmente chi li abbia scoperti per primo. % lang-it

Uczeni arabscy stwierdzą, że był to Archimedes, choć nie ma ku temu niezbitych dowodów. % lang-pl
Gli studiosi arabi affermeranno che fu Archimede, anche se non esistono prove definitive. % lang-it
\index[persons]{Archimede}% % lang-it

\index[persons]{Archimedes}% % lang-pl
Pappus będzie świadom istnienia ortocentrum, ale najstarsze znane uzasadnienie (w komentarzu al-Nasawiego) przyjdzie na świat dopiero w XI wieku. % lang-pl
Pappo conoscerà l'esistenza dell'ortocentro, ma la più antica dimostrazione nota (in un commento di al-Nasawi) apparirà soltanto nell'XI secolo. % lang-it
\index[persons]{Pappus}% % lang-pl
\index[persons]{Pappo}% % lang-it
\index[persons]{al-Nasawi, Ali ibn Ahmad}%

Fakt \ref{srodkowe_przecinaja_sie} ma trójwymiarową kontynuację: % lang-pl
Il fatto \ref{srodkowe_przecinaja_sie} possiede un'analoga versione tridimensionale: % lang-it
\begin{proposition} % lang-it
    I quattro segmenti che uniscono i vertici di un tetraedro con i baricentri delle facce opposte si incontrano in un unico punto, che li divide nel rapporto $3$ a $1$. % lang-it
\end{proposition} % lang-it


\begin{proposition} % lang-pl
    Cztery odcinki łączące wierzchołki czworościanu ze środkami przeciwległych ścian przecinają się w~jednym punkcie, który dzieli je w stosunku $3$ do $1$. % lang-pl
\end{proposition} % lang-pl
Alcuni usano l'espressione ``teorema di Commandino'', poiché Federico Commandino scriverà nel 1565 l'opera \emph{De centro gravitatis solidorum} (Sui centri di gravità dei solidi), pur non essendo necessariamente il primo ad averlo scoperto. % lang-it
\index{teorema!di Commandino}% % lang-it
\index[persons]{Commandino, Federico}% % lang-it

Niektórzy stosują określenie ,,twierdzenie Commandino'', ponieważ Federico Commandino napisze w 1565 roku pracę \emph{De centro gravitatis solidorum} (O środkach ciężkości brył), chociaż może nie być pierwszym, który je odkrył. % lang-pl
Si sospetta che Francesco Maurolico lo conoscesse già in precedenza. % lang-it

\index{twierdzenie!Commandino}% % lang-pl
\index[persons]{Commandino, Federico}% % lang-pl
Podejrzewa się, że Francesco Maurolico pozna je wcześniej. % lang-pl
(Friederich Eduard Reusch znajdzie uogólnienie, które w zdegenerowanym przypadku prowadzi znowu do twierdzenia \ref{theorem_varignon} Varignona). % lang-pl
(Friederich Eduard Reusch troverà una generalizzazione che, in un caso degenere, conduce nuovamente al teorema di Varignon \ref{theorem_varignon}). % lang-it
\index[persons]{Reusch, Friederich Eduard}%

%